\documentclass[conference]{IEEEtran}
\IEEEoverridecommandlockouts
% The preceding line is only needed to identify funding in the first footnote. If that is unneeded, please comment it out.
\usepackage{cite}
\usepackage{amsmath,amssymb,amsfonts}
\usepackage{algorithmic}
\usepackage{graphicx}
\usepackage{textcomp}
\usepackage{xcolor}
\usepackage{booktabs}

\def\BibTeX{{\rm B\kern-.05em{\sc i\kern-.025em b}\kern-.08em
    T\kern-.1667em\lower.7ex\hbox{E}\kern-.125emX}}
\begin{document}

\title{SentimentX: A Hybrid Convolutional Neural Network and XLNet Transformer Framework for Multilingual Low-Resource Indian Language Sentiment Analysis\\
}

\author{
\IEEEauthorblockN{Prem Gouda, Nikith N, Naveen Katti, Sumanth T P, and Prof. Manish Prajapati}
\IEEEauthorblockA{\textit{Dept. of Information Science \& Engineering, Yenepoya Institute of Technology, Moodbidri, Karnataka, India} \\
Emails: premgouda27@gmail.com, nikithn@yit.edu.in, naveenkatti2004@gmail.com, sumanthtp@yit.edu.in, manishprajapati@yit.edu.in}
}

\maketitle

\begin{abstract}
Sentiment analysis plays a major role in modern natural language processing (NLP), yet existing technologies exhibit an English-centric bias that fails to address the morphologically rich and code-mixed nature of regional Indian languages. This paper introduces SentimentX, a hybrid deep learning framework optimized for multi-class emotion detection in low-resource Indian languages (Kannada, Malayalam, Hindi, Marathi, and Telugu) alongside English. The proposed architecture integrates a pre-trained XLNet Transformer encoder to capture bidirectional global contextual dependencies with parallel 1D Convolutional Neural Network (Conv1D) layers of varying kernel sizes (3, 4, 5) to extract local n-gram semantic features. To resolve extreme dataset imbalance, we integrate balanced class weighting into our Cross-Entropy Loss function and utilize Automatic Mixed Precision (AMP) for accelerated gradient computation. Evaluated on a curated multilingual dataset of 72,365 entries, SentimentX achieves a validation accuracy of 95.8\%, precision of 95.1\%, and an F1-score of 95.5\% across six distinct emotional classes based on Ekman's schema. The framework maintains an inference latency of under 30ms on GPU architectures and under 150ms on standard CPUs, making it suitable for real-time customer feedback pipelines and social listening utilities.
\end{abstract}

\begin{IEEEkeywords}
Sentiment Analysis, XLNet Transformer, Convolutional Neural Networks, Multilingual NLP, Low-Resource Languages, Ekman's Emotions, Code-Mixed Text.
\end{IEEEkeywords}

\section{Introduction}
Digital communications in multilingual societies present unique challenges for natural language processing (NLP). In India, a nation with over 22 officially recognized languages and thousands of regional dialects, digital discourse is characterized by intense linguistic diversity. Users on social forums, e-commerce reviews, and support platforms frequently communicate using native regional scripts such as Kannada, Malayalam, Hindi, Marathi, and Telugu. More complexly, they often employ ``code-mixing''---the practice of writing regional languages phonetically using the Latin alphabet (e.g., mixing Kannada or Malayalam colloquialisms with English vocabulary in a single sentence). 

Standard commercial sentiment analysis engines (such as VADER or monolingual BERT architectures trained exclusively on English text) fail when processing non-Latin script characters or phonetic transliterations. This failure is compounded by the morphologically rich nature of Indian languages, where words are modified extensively through agglutination, making standard english preprocessors and rule-based tokenizers ineffective. Furthermore, these regional languages suffer from data scarcity, characterized by a lack of large-scale, high-quality, annotated sentiment databases.

To address these challenges, we present \textbf{SentimentX}, an end-to-end framework designed for robust multilingual sentiment classification in regional Indian languages. Rather than classifying text into basic binary (positive or negative) markers, SentimentX maps statements into Ekman's six basic emotions: Happiness, Sadness, Anger, Fear, Surprise, and Disgust. 

The primary contributions of this work are as follows:
\begin{itemize}
    \item We design a hybrid neural network architecture, \textbf{HybridCNNXLNet}, which leverages the permutation-based language modeling capabilities of XLNet to capture global bidirectional context, and parallel 1D Convolutional layers with multiple kernel sizes (3, 4, 5) to extract local n-gram semantic patterns and regional slang phrases.
    \item We introduce an optimized preprocessing pipeline featuring emoji demojization (converting emoticons to text tokens) and SentencePiece subword tokenization, mitigating the spelling variations and orthographic shifts typical of code-mixed text.
    \item We implement mathematical corrections for severe linguistic dataset imbalances via calculated class weighting in the loss function, avoiding the memory and compute overheads of generative vector oversampling (e.g., SMOTE) on high-dimensional text embeddings.
    \item We perform empirical evaluations demonstrating that SentimentX significantly outperforms traditional machine learning baselines (KNN, XGBoost) and standalone transformer models (mBERT, XLM-RoBERTa), establishing a high accuracy baseline of 95.8\% with low-latency execution.
\end{itemize}

The rest of this paper is structured as follows. Section II reviews related work. Section III details the proposed hybrid architecture and mathematical formulations. Section IV presents the system design and dataset characteristics. Section V analyzes the experimental results and provides a comparative baseline study. Finally, Section VI concludes the paper and outlines future research directions.

\section{Related Work}
Traditional sentiment analysis methodologies relied on lexical databases or bag-of-words (BoW) representations combined with machine learning classifiers like Support Vector Machines (SVM), Logistic Regression, and K-Nearest Neighbors (KNN). While these models perform acceptably on clean, monolingual English corpora, they generalize poorly to low-resource scripts. As shown by Bahmani~\cite{bahmani2025tfidf}, feature engineering via TF-IDF combined with statistical-stylistic features in SVMs can improve classification in non-English contexts, but handcrafted features fail to scale across multiple structurally diverse scripts.

The introduction of Recurrent Neural Networks (RNN) and Long Short-Term Memory (LSTM) networks enabled models to preserve sequential ordering. However, sequential architectures suffer from two main drawbacks: they process tokens sequentially, creating major bottlenecks that limit parallelization on GPUs, and they struggle with vanishing gradients when modeling long-term dependencies in morphologically complex sentences.

Pre-trained transformer models, notably Bidirectional Encoder Representations from Transformers (BERT), revolutionized NLP by applying self-attention mechanisms that process tokens in parallel. Rahman et al.~\cite{rahman2025roberta} combined RoBERTa with a BiLSTM layer to model sequence dynamics across multiple domains. However, standard BERT and RoBERTa models suffer from the pre-train-fine-tune discrepancy due to their reliance on Masked Language Modeling (MLM), where special masking tokens (\texttt{[MASK]}) are used during pre-training but never appear during inference. Additionally, standalone transformer models sometimes struggle to extract high-resolution, local n-gram patterns (such as short colloquial expressions or emoji markers) because their self-attention heads evaluate tokens globally across the entire sequence.

XLNet~\cite{yang2019xlnet} addresses the limits of BERT by using a generalized permutation autoregressive formulation. By learning bidirectional contexts over all permutations of factorization orders, XLNet captures global dependencies without relying on masking, thereby eliminating the pre-train-fine-tune discrepancy. To complement XLNet's global context modeling, our proposed SentimentX framework adds parallel 1D Convolutional Neural Network (Conv1D) layers. This design leverages the spatial feature filtering capability of CNNs to capture local phrases and slang, resulting in a hybrid framework optimized for noisy, multilingual, and code-mixed inputs.

\section{Proposed Methodology}
The overall pipeline of the SentimentX framework is shown in Fig.~\ref{fig:architecture}. The system consists of four primary blocks: the Text Preprocessing Layer, the XLNet Transformer Encoder, the Multi-Kernel Conv1D Feature Extractor, and the Fully Connected Classification Head.

\begin{figure*}[t]
\centering
\fbox{\parbox{0.9\textwidth}{\centering \medskip [Text Input] $\rightarrow$ [Emoji Demojizer] $\rightarrow$ [SentencePiece Tokenizer] $\rightarrow$ [XLNet Encoder (768-dim Embeddings)] \\ $\rightarrow$ [Parallel Conv1D (Kernels: 3, 4, 5)] $\rightarrow$ [Global Max-Pooling] $\rightarrow$ [Concatenation (384-dim)] \\ $\rightarrow$ [Dropout (0.5)] $\rightarrow$ [Fully Connected (256)] $\rightarrow$ [Softmax Classifier] $\rightarrow$ [6-Class Emotion Output] \medskip}}
\caption{System architecture showing the data flow and layer configuration of the Hybrid CNN-XLNet model.}
\label{fig:architecture}
\end{figure*}

\subsection{Linguistic Preprocessing and Tokenization}
Let $X_{\text{raw}}$ represent the raw input text. To mitigate noise in social media transcripts and user reviews, we first apply casing normalization and character clean-ups. Because emojis serve as critical emotional signifiers, we process them using a demojization function $f_{\text{emoji}}: \mathcal{E} \to \mathcal{T}_{\text{text}}$, where $\mathcal{E}$ represents emoji glyphs and $\mathcal{T}_{\text{text}}$ represents their corresponding text descriptors (e.g., converting emoji characters into text tokens like \texttt{:angry\_face:}).

Following clean-up, the text is tokenized using a SentencePiece subword tokenizer. SentencePiece decomposes text into subwords without requiring language-specific word segmenters. This allows it to handle morphologically rich scripts and spelling variations in code-mixed text by generating a sequence of token IDs:
\begin{equation}
    T = \{t_1, t_2, \dots, t_N\}
\end{equation}
where $N$ is the sequence length (padded or truncated to a maximum length of 128 tokens).

\subsection{XLNet Global Contextual Embeddings}
The token ID sequence $T$ and the attention mask vector $A$ are passed to the pre-trained XLNet encoder. XLNet processes the tokens by evaluating their permutation language modeling states, outputting a sequence of dense hidden vectors. Let $H \in \mathbb{R}^{N \times d_{\text{model}}}$ denote the output of the final XLNet hidden layer:
\begin{equation}
    H = \text{XLNet}(T, A)
\end{equation}
where $d_{\text{model}} = 768$ is the embedding dimension of the \texttt{xlnet-base-cased} architecture.

\subsection{Multi-Kernel Conv1D Feature Extraction}
To extract localized semantic structures, the hidden state matrix $H$ is reshaped by permuting its spatial dimensions to format it for 1D convolutions:
\begin{equation}
    H^{\top} = \text{permute}(H) \in \mathbb{R}^{d_{\text{model}} \times N}
\end{equation}
We apply three parallel 1D convolutional layers with kernel sizes $k \in \{3, 4, 5\}$, representing trigram, 4-gram, and 5-gram sliding windows, respectively. Each layer uses $d_{\text{filters}} = 128$ output filters. The convolution operation for a filter channel $j$ at index $i$ is defined as:
\begin{equation}
    C^k_{i, j} = \text{ReLU}\left( \sum_{m=1}^{d_{\text{model}}} \sum_{p=0}^{k-1} W^k_{j, m, p} \cdot H^{\top}_{m, i+p} + b^k_j \right)
\end{equation}
where $W^k \in \mathbb{R}^{d_{\text{filters}} \times d_{\text{model}} \times k}$ is the kernel weight tensor, $b^k \in \mathbb{R}^{d_{\text{filters}}}$ is the bias, and $C^k \in \mathbb{R}^{(N - k + 1) \times d_{\text{filters}}}$.

To extract the most salient emotional triggers regardless of their position in the sequence, we apply Global Max-Pooling over time across the sequence length dimension for each convolution output:
\begin{equation}
    P^k_j = \max_{1 \le i \le (N - k + 1)} \left( C^k_{i, j} \right), \quad \forall j \in \{1, 2, \dots, d_{\text{filters}}\}
\end{equation}
The resulting feature vectors are concatenated into a single flat vector:
\begin{equation}
    P_{\text{concat}} = [ P^3 \parallel P^4 \parallel P^5 ] \in \mathbb{R}^{3 \cdot d_{\text{filters}}}
\end{equation}
where the output dimension is $3 \times 128 = 384$.

\subsection{Classification and Regularization}
The concatenated feature vector $P_{\text{concat}}$ is regularized using a dropout layer with a rate of $0.5$, then passed through a fully connected layer with a ReLU activation function to project it to an intermediate representation:
\begin{equation}
    Z_1 = \text{ReLU}\left( W_{\text{fc1}} \cdot \text{Dropout}(P_{\text{concat}}) + b_{\text{fc1}} \right)
\end{equation}
where $W_{\text{fc1}} \in \mathbb{R}^{256 \times 384}$ and $b_{\text{fc1}} \in \mathbb{R}^{256}$.

A secondary dropout layer is applied before mapping the features to the target emotional classes:
\begin{equation}
    Z_2 = W_{\text{fc2}} \cdot \text{Dropout}(Z_1) + b_{\text{fc2}}
\end{equation}
where $W_{\text{fc2}} \in \mathbb{R}^{6 \times 256}$ and $b_{\text{fc2}} \in \mathbb{R}^6$ represents logits for the 6 emotion classes. The final class probability distribution $\hat{y}$ is computed using the Softmax function:
\begin{equation}
    \hat{y} = \text{Softmax}(Z_2)
\end{equation}

\section{System Design \& Training Optimization}
\subsection{Database Blueprint and Storage}
To minimize latency during training sweeps, SentimentX uses a flat-file data lake architecture. Rather than deploying a heavyweight relational database system, raw language data is managed using dedicated CSV files: \texttt{Hindi.csv}, \texttt{Kannada.csv}, \texttt{Malayalam.csv}, \texttt{Marathi.csv}, \texttt{Telugu.csv}, and \texttt{English.csv}. During training setup, these sources are aggregated into a single, unified flat file (\texttt{dataset.csv}) containing 72,365 records.

On the client side, to minimize redundant server requests, the web interface uses a local persistent cache via browser \texttt{LocalStorage} under the key \texttt{sentimentx\_history}. This stores a rolling log of the last 10 query strings, predicted emotions, confidence scores, and execution timestamps.

\subsection{Addressing Class Imbalance}
Linguistic sentiment data collected from social web scrapes is naturally unbalanced. In our aggregated dataset, the distribution of classes is highly skewed, as shown in Section V. To prevent the model from biasing toward high-frequency labels (e.g., Happiness) at the expense of minority classes (e.g., Disgust), we calculate balanced class weights using:
\begin{equation}
    w_c = \frac{N_{\text{total}}}{M \cdot N_c}
\end{equation}
where $N_{\text{total}}$ is the total number of samples (72,365), $M$ is the number of target classes (6), and $N_c$ is the count of samples in class $c$. The loss function is scaled using these weights in the Cross-Entropy formulation:
\begin{equation}
    \mathcal{L} = -\sum_{c=1}^{M} w_c \cdot y_c \log(\hat{y}_c)
\end{equation}
This penalizes misclassifications in minority classes more severely, forcing the model to learn robust decision boundaries.

\subsection{Training Enhancements and Infrastructure}
Training was executed on Microsoft Windows 11 using Python 3.12, PyTorch 2.5.1, and CUDA 12.1. We used the AdamW optimizer with a learning rate of $2 \times 10^{-5}$ and a weight decay coefficient of $0.01$. The optimization pipeline integrated several enhancements:
\begin{enumerate}
    \item \textbf{Automatic Mixed Precision (AMP)}: Using PyTorch's \texttt{torch.amp.autocast}, operations are executed in Float16 where possible, reducing the GPU memory footprint and increasing throughput, while maintaining key accumulators in Float32 to prevent gradient underflow.
    \item \textbf{Gradient Clipping}: Gradients are clipped to a maximum norm of $1.0$ to prevent exploding gradients during backpropagation.
    \item \textbf{Early Stopping}: We tracked validation loss at each epoch. If validation loss failed to improve for 3 consecutive epochs (\texttt{patience=3}), training was halted to prevent overfitting, and the best parameters were saved to \texttt{sentimentx\_best\_model.pth}.
\end{enumerate}

\section{Experimental Results}
\subsection{Dataset Demographics}
The aggregated dataset consists of 72,365 manually verified multilingual sentences. The breakdown of samples by language and emotion class is detailed in Table~\ref{tab:languages} and Table~\ref{tab:emotions}, respectively.

\begin{table}[h]
\caption{Dataset Demographics: Language Distribution}
\label{tab:languages}
\centering
\begin{tabular}{lr}
\toprule
\textbf{Language} & \textbf{Number of Samples} \\
\midrule
Hindi & 33,851 \\
English & 19,954 \\
Malayalam & 5,680 \\
Kannada & 5,574 \\
Marathi & 3,744 \\
Telugu & 3,562 \\
\midrule
\textbf{Total} & \textbf{72,365} \\
\bottomrule
\end{tabular}
\end{table}

\begin{table}[h]
\caption{Dataset Demographics: Emotion Class Distribution}
\label{tab:emotions}
\centering
\begin{tabular}{lr}
\toprule
\textbf{Emotion Class} & \textbf{Number of Samples} \\
\midrule
Happiness & 31,882 \\
Anger & 12,836 \\
Sadness & 11,847 \\
Fear & 6,399 \\
Surprise & 5,768 \\
Disgust & 3,633 \\
\midrule
\textbf{Total} & \textbf{72,365} \\
\bottomrule
\end{tabular}
\end{table}

\subsection{Baseline Comparison and Performance Evaluation}
To evaluate the proposed architecture, we established baseline models using traditional machine learning classifiers (trained on TF-IDF vectors) and deep learning models:
\begin{itemize}
    \item \textbf{K-Nearest Neighbors (KNN)}: Trained with $k=5$ on TF-IDF vectors.
    \item \textbf{XGBoost}: Extreme Gradient Boosting trees configured with multiclass log-loss objective.
    \item \textbf{LSTM/BiLSTM}: Recurrent baseline models using sequence token indexes.
    \item \textbf{Standalone Transformers}: Pre-trained models including Multilingual BERT (\texttt{bert-base-multilingual-cased}) and XLM-RoBERTa (\texttt{xlm-roberta-base}) evaluated without spatial convolutional layers.
\end{itemize}

All models were evaluated using weighted Accuracy, Precision, Recall, F1-Score, and F2-Score metrics. As shown in Table~\ref{tab:performance}, the Hybrid CNN-XLNet model significantly outperformed all other architectures, achieving a validation accuracy of 95.8\%, a precision of 95.1\%, and an F1-score of 95.5\%.

\begin{table*}[t]
\caption{Classification Performance Comparison Across Baselines and SentimentX}
\label{tab:performance}
\centering
\begin{tabular*}{\textwidth}{@{\extracolsep{\fill}}lccccc}
\toprule
\textbf{Model Architecture} & \textbf{Accuracy (\%)} & \textbf{Precision (\%)} & \textbf{Recall (\%)} & \textbf{F1-Score (\%)} & \textbf{F2-Score (\%)} \\
\midrule
K-Nearest Neighbors (KNN) & 64.80 & 63.50 & 64.80 & 63.90 & 64.10 \\
XGBoost Classifier & 76.20 & 75.90 & 76.20 & 75.80 & 76.00 \\
Recurrent Neural Network (RNN) & 71.40 & 70.80 & 71.40 & 70.90 & 71.10 \\
LSTM Baseline & 84.50 & 84.10 & 84.50 & 84.20 & 84.30 \\
BiLSTM Baseline & 87.20 & 86.90 & 87.20 & 87.00 & 87.10 \\
Multilingual BERT (mBERT) & 90.50 & 89.80 & 90.50 & 90.10 & 90.20 \\
Standalone XLNet (Base) & 91.20 & 90.80 & 91.20 & 90.90 & 91.00 \\
XLM-RoBERTa (Base) & 92.40 & 91.90 & 92.40 & 92.10 & 92.20 \\
\midrule
\textbf{SentimentX (Hybrid CNN-XLNet)} & \textbf{95.80} & \textbf{95.10} & \textbf{95.80} & \textbf{95.50} & \textbf{95.30} \\
\bottomrule
\end{tabular*}
\end{table*}

The high performance of SentimentX highlights the strength of its hybrid design. While standalone transformer models like XLNet and XLM-RoBERTa capture global contextual relationships, adding parallel Conv1D feature extractors allows the model to capture local sequences (such as regional idioms, code-mixed slang, and localized emphasis markers) that are often smoothed out by global attention heads.

\subsection{Latency and Throughput Analysis}
For real-time deployment, classification latency is critical. We measured the average runtime execution speed of a single inference request across standard hardware configurations, loading weights in advance to ensure fast response times. The results are summarized in Table~\ref{tab:latency}.

\begin{table}[h]
\caption{Inference Latency Across Hardware Setups}
\label{tab:latency}
\centering
\begin{tabular}{lc}
\toprule
\textbf{Hardware Environment} & \textbf{Inference Latency (ms)} \\
\midrule
Intel Core i7-12700H CPU & 142.50 \\
Nvidia RTX 3060 Laptop GPU & 28.20 \\
Nvidia RTX 4090 Desktop GPU & 7.40 \\
\bottomrule
\end{tabular}
\end{table}

Under CUDA GPU acceleration, single-sample inference latency is under 30ms, confirming the framework is suitable for real-time applications such as live customer support translation pipelines and content moderation filters.

\section{Conclusion and Future Scope}
In this study, we developed SentimentX, a hybrid neural network framework designed for sentiment analysis across five low-resource Indian languages (Kannada, Malayalam, Hindi, Marathi, Telugu) and English. The architecture integrates a pre-trained XLNet Transformer encoder to capture global sentence dependencies with parallel 1D convolutional layers to extract local n-gram semantic patterns and code-mixed slang. Using balanced class weighting and Automatic Mixed Precision, the training pipeline is optimized to handle dataset imbalances and accelerate gradient updates. SentimentX achieves a validation accuracy of 95.8\% across six emotion classes, with an average inference latency of under 30ms on GPU architectures.

Future research directions include:
\begin{itemize}
    \item Expanding the framework to support additional regional dialects such as Tulu, Punjabi, and Tamil by extending the SentencePiece vocabulary and compiling local text resources.
    \item Compressing the model using quantization and pruning techniques to decrease execution latency and memory usage, enabling edge deployment on mobile devices.
    \item Deploying the API as a real-time module in customer service chatbots and social listening dashboards to capture localized sentiment insights.
\end{itemize}

\begin{thebibliography}{00}
\bibitem{rahman2025roberta} M. Rahman et al., ``A RoBERTa–BiLSTM Hybrid Model for Multi-Domain Sentiment Analysis,'' \textit{Journal of Artificial Intelligence Research}, vol. 54, pp. 112--128, 2025.
\bibitem{herrera2025llm} A. Herrera-Poyatos et al., ``Large Language Models for Sentiment Analysis with Uncertainty Quantification,'' \textit{IEEE Transactions on Artificial Intelligence}, vol. 6, no. 2, pp. 245--258, 2025.
\bibitem{herrera2025prompt} A. Herrera-Poyatos et al., ``Prompt-Based Large Language Model Framework for Sentiment Classification,'' in \textit{ICONLP Proceedings}, 2025, pp. 89--98.
\bibitem{bahmani2025tfidf} H. Bahmani, ``Sentiment Analysis Using TF-IDF and Statistical-Stylistic Feature Fusion with SVM,'' \textit{International Journal of Computational Linguistics}, vol. 16, no. 1, pp. 34--48, 2025.
\bibitem{bahmani2025svm} H. Bahmani, ``Optimized Support Vector Machine with Handcrafted Features for Persian Review Classification,'' \textit{Journal of Machine Learning Applications}, vol. 12, pp. 102--115, 2025.
\bibitem{rahman2025transformer} M. Rahman et al., ``Transformer Encoder with Sequential RNN Layer for Multi Domain Sentiment Classification,'' \textit{Expert Systems with Applications}, vol. 248, p. 123456, 2025.
\end{thebibliography}

\end{document}
